\chapter{Experimental Setup}
\label{ch:experimental_setup}

\todoX{Write a brief introduction summarizing the purpose of this chapter: to define the hardware, software, evaluation metrics, and exact physical and algorithmic parameters required to reproduce the global trajectory optimization benchmarks.}

\section{Computational Environment \& Software Stack}
\label{sec:setup_environment}
\todoX{Detail hardware specifications (e.g., CPU, RAM) and operating system.}
\todoX{Specify framework architecture (e.g., Python orchestration interacting with a C++ core).}
\todoX{Explicitly list libraries used: CasADi for AD, IPOPT for the NLP, and specific linear solvers (e.g., MUMPS, HSL MA57).}
\todoX{State universal solver stopping criteria (IPOPT convergence tolerance $\varepsilon_{tol} = 10^{-6}$, max iterations).}

The experiments were conducted on two workstations running Ubuntu 24.04.5 LTS on an x86\textsubscript{64} architecture. 
The first computer is equipped with an Intel\textsuperscript{\textregistered} Xeon\textsuperscript{\textregistered} CPU E5-1630 v4 @ 3.70\,GHz, providing 8 logical processors and 16\,GB of RAM. 
The second computer is equipped with an Intel\textsuperscript{\textregistered} Core\textsuperscript{\texttrademark} Ultra 7 265 processor, providing 20 logical processors and 32\,GB of RAM.
As these are materially different in terms of CPU and RAM, the performance results will be reported with with notes indicating which workstation was used for each experiment.

The code framework is based on a Python orchestration layer that interacts with a C++ core for computational efficiency. 
\ac{CasADi} constructs the symbolic B-spline dynamics, objective, constraints, and automatic differentiation, while \ac{IPOPT} solves the resulting nonlinear programs with HSL MA27 as the linear solver. 
NumPy provides vectorized numerical operations, including trajectory sampling. 
SciPy provides selected interpolation, B-spline, spatial, and rotation utilities used by parts of the framework. 

\ac{IPOPT} convergence criteria are unique to each solver as shown in Table~\ref{tab:solver-stopping-criteria}. All other unspecified stopping criteria use the \ac{IPOPT} defaults:

\begin{table}[H]
	\centering
	\scriptsize
	\caption{Solver-specific stopping criteria used in the experiments.}
	\label{tab:solver-stopping-criteria}
		\begin{tabular}{@{}>{\raggedright\arraybackslash}p{0.13\textwidth}>{\raggedright\arraybackslash}p{0.13\textwidth}>{\raggedright\arraybackslash}p{0.10\textwidth}>{\raggedright\arraybackslash}p{0.15\textwidth}>{\raggedright\arraybackslash}p{0.34\textwidth}@{}}
		\toprule
			\makecell[l]{\textbf{Algorithm}} & \makecell[l]{\textbf{Solver stage}} & \makecell[l]{\textbf{IPOPT}\\\textbf{max\_iter}} & \makecell[l]{\textbf{IPOPT}\\\textbf{tol}} & \makecell[l]{\textbf{Other stopping}\\\textbf{criteria}} \\
		\midrule
		MC, both systems & Single-stage solve & 500 & $10^{-10}$ & \parbox[t]{\linewidth}{\raggedright\texttt{constr\_viol\_tol=1e-4}} \\
		P-HO, both systems & Intermediate solve & 500 & $10^{-2}$ & \parbox[t]{\linewidth}{\raggedright\texttt{acceptable\_tol=5e-2}; \\ \texttt{acceptable\_iter=3}; \\ \texttt{constr\_viol\_tol=1e-2}} \\
		P-HO, both systems & Final solve & 500 & $10^{-10}$ & \parbox[t]{\linewidth}{\raggedright\texttt{constr\_viol\_tol=1e-4}} \\
		\parbox[t]{\linewidth}{\raggedright RRT$^{*}$-GBO, \\ both systems} & Intermediate edge solve & 500 & $10^{-2}$ & \parbox[t]{\linewidth}{\raggedright\texttt{acceptable\_tol=5e-2}; \\ \texttt{acceptable\_iter=3}; \\ \texttt{constr\_viol\_tol=1e-2}} \\
		\parbox[t]{\linewidth}{\raggedright RRT$^{*}$-GBO, \\ both systems} & Final solve & 500 & $10^{-10}$ & \parbox[t]{\linewidth}{\raggedright\texttt{constr\_viol\_tol=1e-4}} \\
		IPSO, spacecraft & IPSO stage & N/A & N/A & \parbox[t]{\linewidth}{\raggedright\texttt{convergence\_window=50}; \\ \texttt{convergence\_tolerance}=$1.0\times10^{-6}$; \\ Eq.~\eqref{eq:ipso-convergence}} \\
		IPSO, spacecraft & Warmstarted GBO stage & 500 & $10^{-10}$ & \parbox[t]{\linewidth}{\raggedright\texttt{constr\_viol\_tol=1e-4}} \\
		\bottomrule
	\end{tabular}
\end{table}

\section{Evaluation Metrics}
\label{sec:setup_metrics}
\todoX{Define Success Rate: e.g., IPOPT returns "Solve Succeeded" and the adaptive refinement detects zero continuous-time violations.}
\todoX{Define Computational Effort: Wall-clock CPU time, including notes on parallelization/multi-threading for MC-GBO and IPSO.}
\todoX{Define Solution Quality: The final objective cost J achieved (effort for spacecraft, effort + acceleration for the manipulator).}

The evaluation is organized into four metric groups. 
Reliability and computation time are reported for all attempts, whereas solution quality is evaluated only for successful solutions. 
Search-behavior metrics are reported for repeated runs and for the algorithms for which the corresponding internal statistics are available. 
The individual results sections specify the applicable metrics and reporting subsets for each experiment.

\paragraph{Reliability}
Reliability is characterized by the success rate and failure rate over all attempted runs. 
\ac{IPOPT} failures are classified as infeasible problem instances, maximum-iteration termination, timeouts, or other unexpected exceptions. 
A run is successful only when the solver returns converges with the resulting trajectory passing the applicable kinodynamic constraint checks.

\paragraph{Computation Time}
Computation time includes per test case statistics on wall-clock time, sequential-equivalent time for parallelized algorithms, time to first solution, and time to best solution. 
Solve-time distributions are summarized using the median, mean, standard deviation, and relevant quantiles. 
Quantile solve durations are reported where the number of trials supports a meaningful distributional comparison.

\paragraph{Solution Quality}
For successful solutions, the final objective cost is reported using the mean, standard deviation, median, and relevant quantiles, together with the best, worst, and interquartile-range costs when appropriate. 
Results also include the relative cost gap to the best solution found across all algorithms. 
The objective definition is problem-specific: spacecraft results use the control effort objective, while manipulator results include the specified control effort plus acceleration terms.

\paragraph{Search Behavior and Diversity}
Search behavior is analyzed to characterize local-minimum discovery and solution diversity. 
Metrics include the raw number of solutions, unique basin frequencies. 
Algorithm-specific statistics are reported when available: for \ac{P-HO}, final $\lambda=1$ solutions, tree size; for \ac{RRT*-GBO}, sampled nodes, accepted edges; and for \ac{IPSO}, particle evaluations, completed outer and inner iterations, and the number of solutions sent to \ac{GBO} refinement.

The complete metric inventory is defined here for consistency, but not every metric is relevant to every experiment. 
The results section therefore reports the metrics applicable to the algorithm and test case under study, while retaining the common definitions above for cross-algorithm comparisons.

\section{Problem 1: Spacecraft Attitude Slew}
\label{sec:setup_spacecraft}

\subsection{Physical Parameters}
\todoX{Define the spacecraft inertia matrix $\mathbf{J}_{\mathrm{sc}}$.}
\todoX{List actuator torque limits $\boldsymbol{\tau}_{\max}$ and angular velocity limits $\boldsymbol{\omega}_{\max}$.}
The spacecraft inertia matrix denoted by $\mathbf{J}_{\mathrm{sc}}$ is set as a spherical inertia, $\mathbf{J}_{\mathrm{sc}} = 10000*\mathbf{I}$ kg·m².
The baseline control torque limits are set to $\boldsymbol{\tau}_{\max} = 200$ N·m, and the angular velocity limits are set to $\boldsymbol{\omega}_{\max} = 0.2$ rad/s (11.46 °/s).
If different torque or angular velocity limits are used for specific test cases, they are explicitly stated in the scenario description.
Both limits are defined in the spacecraft body frame.

\subsection{B-Spline Discretization}
\todoX{Specify the fixed time horizon $t_f$.}
\todoX{Specify B-spline degree (e.g., $p=3$ or $p=4$) and number of control vertices $n_{ctrl}$.}
\todoX{State the oversampling ratio ($n_{via} \approx 4 n_{ctrl}$).}
For all test scenarios, the B-splines were parameterized via:
\begin{itemize}
    \item Fixed time horizon: $t_f = 60$ s
    \item B-spline degree: cubic $p = 3$
    \item Number of free control vertices: $n_{FCV} = 10$
    \item Number of boundary conditions: $n_{BC} = 3^\dagger$ equating to a set angular orientation, zero angular velocity, and zero angular acceleration, $^\dagger$ except for two test cases that include additionally zero angular jerk $n_{BC} = 4$.
    \item Total number of control vertices: $n_{ctrl} = n_{FCV} + 2n_{BC}$
    \item Number of via points: $n_{via} = 61$
\end{itemize}

\subsection{Test Scenarios}
\todoX{Define specific boundary conditions ($\mathbf{q}_0$, $\mathbf{q}_f$, etc.).}
\todoX{Case A: Unconstrained large-angle slew (baseline).}
\todoX{Case B: Constrained slew with Exclusion Zones (EZs). Provide a table or list of inertial pointing vectors ${}^i\mathbf{\hat{e}}$ and radii $\theta_{exclusion}$.}

Tables~\ref{tab:scso3_summary} and \ref{tab:scso3_exclusion_zones} summarize the spacecraft slew test case scenarios and the corresponding exclusion-zone definitions, respectively.
The test case scenarios range from an unconstrained large-angle slew---as a baseline verification of the optimization framework---to increasingly constrained realistic maneuvers with exclusion zones and actuator limits, and finally complex configuration scenarios testing the limits of the global solver's ability to find feasible and optimal trajectories. 

In addition to the listed orientation vectors of the \ac{EZ}s, the corresponding sensor vector is defined as ${}^b\mathbf{\hat{s}} = \left[0, 0, 1\right]^{T}$, which is transformed to the inertial frame via the current attitude of the spacecraft for comparison with the exclusion-zone directions.

\newcommand{\quat}[4]{\begin{gathered}\text{[}#1,\ #2,\\[-1pt]#3,\ #4\text{]}^{T}\end{gathered}}
\newcommand{\ezvec}[3]{\left[#1,\ #2,\ #3\right]^{T}}

% ============================================================
% Table: Scenario summary
% ============================================================

\begin{table}[H]
\centering
\scriptsize
\setlength{\tabcolsep}{4pt}
\renewcommand{\arraystretch}{1.2}
\caption{Summary of the spacecraft slew test case scenarios.}
\label{tab:scso3_summary}
\begin{tabularx}{\textwidth}{
    >{\raggedright\arraybackslash}p{3.2cm}
	>{\centering\arraybackslash}p{3.8cm}
	>{\centering\arraybackslash}p{3.8cm}
    >{\centering\arraybackslash}p{0.8cm}
    >{\raggedright\arraybackslash}X
}
\toprule
\textbf{Test case}
& $\mathbf q_0$
& $\mathbf q_f$
& \textbf{\# EZs}
& \textbf{Constraints and boundary conditions} \\
\midrule

\texttt{geodesic slew} & $\quat{1}{0}{0}{0}$ & $\quat{0}{0.563122}{0.563122}{-0.604803}$ & 0 & EZ off; $\tau_{\max}$ off; $\omega_{\max}$ off. \\

\texttt{sun} & $\quat{0}{0}{0}{1}$ & $\quat{0.7044}{-0.0616}{0.0616}{0.7044}$ & 1 & EZ on; $\tau_{\max}$ on; $\omega_{\max}$ on. \\

\texttt{sun and moon} & $\quat{-0.2897}{0.5727}{0.3462}{0.6843}$ & $\quat{-0.2664}{0.0290}{-0.3852}{0.8831}$ & 2 & EZ on; $\tau_{\max}$ on; $\omega_{\max}$ on. \\

\texttt{sun, moon, and earth 1} & $\quat{0.5577}{0.3861}{0.0545}{0.7328}$ & $\quat{-0.2520}{0.2684}{-0.9272}{-0.0689}$ & 3 & EZ on; $\tau_{\max}$ on; $\omega_{\max}$ on. \\

\texttt{sun, moon, and earth 2} & $\quat{0.5577}{0.3861}{0.0545}{0.7328}$ & $\quat{-0.2520}{0.2684}{-0.9272}{-0.0689}$ & 3 & EZ on; $\tau_{\max}$ on; $\omega_{\max}$ on. \\

\texttt{narrow corridor} & $\quat{0.4670}{0.5309}{0.5225}{0.4764}$ & $\quat{-0.2951}{0.2018}{-0.5403}{0.7618}$ & 3 & EZ on; $\tau_{\max}$ on; $\omega_{\max}$ on; jerk BC on. \\

\texttt{narrow corridor} $+ \omega_{\max}$ & $\quat{0.4670}{0.5309}{0.5225}{0.4764}$ & $\quat{-0.2951}{0.2018}{-0.5403}{0.7618}$ & 3 & EZ on; $\tau_{\max}$ on; $\omega_{\max}=0.08~\mathrm{rad/s}$; jerk BC on. \\

\texttt{bug trap} & $\quat{-0.2340}{0.4977}{0.3554}{0.7558}$ & $\quat{0.3318}{0.7358}{0.5824}{-0.0968}$ & 5 & EZ on; $\tau_{\max}$ on; $\omega_{\max}$ on. \\

\texttt{reverse bug trap} & $\quat{0.3318}{0.7358}{0.5824}{-0.0968}$ & $\quat{-0.2340}{0.4977}{0.3554}{0.7558}$ & 5 & EZ on; $\tau_{\max}$ on; $\omega_{\max}$ on. \\

\texttt{double bug trap} & $\quat{-0.7420}{0.4052}{0.4687}{0.2559}$ & $\quat{0.3823}{0.3889}{0.8366}{-0.0526}$ & 10 & EZ on; $\tau_{\max}$ on; $\omega_{\max}$ on. \\

\texttt{maze} & $\quat{-0.2512}{0.5018}{0.3705}{0.7402}$ & $\quat{-0.5295}{0.6263}{0.5720}{0.0135}$ & 8 & EZ on; $\tau_{\max}$ on; $\omega_{\max}$ on. \\

\bottomrule
\end{tabularx}
\end{table}

% ============================================================
% Table: Exclusion-zone definitions
% ============================================================

\begin{longtable}{
    >{\raggedright\arraybackslash}p{3.4cm}
    >{\centering\arraybackslash}p{0.45cm}
    >{\centering\arraybackslash}p{7.0cm}
    >{\centering\arraybackslash}p{1.4cm}
}
\caption{Exclusion-zone definitions for the spacecraft slew test case scenarios.}
\label{tab:scso3_exclusion_zones}\\
\toprule
\textbf{Test case} & $k$ &
\textbf{Inertial exclusion vector ${}^{i}\hat{\mathbf e}_k$} &
\textbf{$\theta_k$} \\
\midrule
\endfirsthead

\toprule
\textbf{Test case} & $k$ &
\textbf{Inertial exclusion vector ${}^{i}\hat{\mathbf e}_k$} &
\textbf{$\theta_k$} \\
\midrule
\endhead

\bottomrule
\endfoot

\texttt{geodesic slew}
& 0 & - & - \\

\texttt{sun}
& 1 & $\ezvec{-0.1}{-0.5}{0.7}$ & $25^\circ$ \\

\texttt{sun and moon}
& 1 & $\ezvec{0.8605}{0.1795}{0.4767}$ & $30^\circ$ \\
& 2 & $\ezvec{0.3743}{0.7954}{0.4767}$ & $15^\circ$ \\

\texttt{sun, moon and earth 1}
& 1 & $\ezvec{0.2796}{-0.8605}{0.4258}$ & $25^\circ$ \\
& 2 & $\ezvec{0.7754}{-0.4615}{0.4310}$ & $15^\circ$ \\
& 3 & $\ezvec{-0.5176}{0.7392}{-0.4310}$ & $135^\circ$ \\

\texttt{sun, moon and earth 2}
& 1 & $\ezvec{0.8074}{-0.4085}{0.4258}$ & $15^\circ$ \\
& 2 & $\ezvec{0.2796}{-0.8605}{0.4258}$ & $25^\circ$ \\
& 3 & $\ezvec{-0.5176}{0.7392}{-0.4310}$ & $135^\circ$ \\

\texttt{narrow corridor}
& 1 & $\ezvec{0.7954}{0.6061}{0.0000}$ & $25^\circ$ \\
& 2 & $\ezvec{0.7530}{-0.4140}{0.5114}$ & $40^\circ$ \\
& 3 & $\ezvec{0.6790}{0.3937}{0.6196}$ & $10^\circ$ \\

\texttt{narrow corridor} $+ \omega_{\max}$
& 1 & $\ezvec{0.7954}{0.6061}{0.0000}$ & $25^\circ$ \\
& 2 & $\ezvec{0.7530}{-0.4140}{0.5114}$ & $40^\circ$ \\
& 3 & $\ezvec{0.6790}{0.3937}{0.6196}$ & $10^\circ$ \\

\texttt{bug trap}
& 1 & $\ezvec{0.8577}{0.5141}{0.0000}$ & $25^\circ$ \\
& 2 & $\ezvec{0.7370}{0.2236}{0.6378}$ & $25^\circ$ \\
& 3 & $\ezvec{0.1737}{0.6806}{0.7118}$ & $25^\circ$ \\
& 4 & $\ezvec{0.2243}{0.8786}{0.4215}$ & $15^\circ$ \\
& 5 & $\ezvec{0.5775}{0.8162}{0.0201}$ & $15^\circ$ \\

\texttt{reverse bug trap}
& 1 & $\ezvec{0.8577}{0.5141}{0.0000}$ & $25^\circ$ \\
& 2 & $\ezvec{0.7370}{0.2236}{0.6378}$ & $25^\circ$ \\
& 3 & $\ezvec{0.1737}{0.6806}{0.7118}$ & $25^\circ$ \\
& 4 & $\ezvec{0.2243}{0.8786}{0.4215}$ & $15^\circ$ \\
& 5 & $\ezvec{0.5775}{0.8162}{0.0201}$ & $15^\circ$ \\

\texttt{double bug trap}
& 1 & $\ezvec{0.8577}{0.5141}{0.0000}$ & $25^\circ$ \\
& 2 & $\ezvec{0.7370}{0.2236}{0.6378}$ & $25^\circ$ \\
& 3 & $\ezvec{0.1737}{0.6806}{0.7118}$ & $25^\circ$ \\
& 4 & $\ezvec{0.2243}{0.8786}{0.4215}$ & $15^\circ$ \\
& 5 & $\ezvec{0.5775}{0.8162}{0.0201}$ & $15^\circ$ \\
& 6 & $\ezvec{-0.5156}{0.8517}{0.0936}$ & $25^\circ$ \\
& 7 & $\ezvec{-0.2944}{0.5012}{-0.8137}$ & $25^\circ$ \\
& 8 & $\ezvec{-0.8422}{0.3938}{-0.3683}$ & $25^\circ$ \\
& 9 & $\ezvec{-0.1110}{0.7969}{-0.5938}$ & $15^\circ$ \\
& 10 & $\ezvec{-0.1682}{0.9819}{-0.0870}$ & $15^\circ$ \\

\texttt{maze}
& 1 & $\ezvec{0.8577}{0.5141}{0.0000}$ & $25^\circ$ \\
& 2 & $\ezvec{0.7370}{0.2236}{0.6378}$ & $25^\circ$ \\
& 3 & $\ezvec{0.1737}{0.6806}{0.7118}$ & $25^\circ$ \\
& 4 & $\ezvec{0.2243}{0.8786}{0.4215}$ & $15^\circ$ \\
& 5 & $\ezvec{0.5775}{0.8162}{0.0201}$ & $15^\circ$ \\
& 6 & $\ezvec{-0.0827}{0.9722}{-0.2191}$ & $25^\circ$ \\
& 7 & $\ezvec{0.4316}{0.4720}{-0.7687}$ & $25^\circ$ \\
& 8 & $\ezvec{-0.6277}{0.6802}{0.3785}$ & $25^\circ$ \\

\end{longtable}

\subsection{Algorithm-Specific Hyperparameters}
\todoX{introduce the tables}

Tables~\cref{tab:params_mc_gbo,tab:params_pho,tab:params_rrt_gbo,tab:params_ipso} list the comprehensive set of algorithm-specific hyperparameters used in this work for the spacecraft \ac{NLPn}.
Refer back to \cref{sec:algorithm_frameworks} and the papers by Stoneman and Lampariello \cite{Stoneman14}, Pardis et al. \cite{Pardis24}, and Spiller et al. \cite{Spiller16,Spiller18} for further details on the respective algorithms.

% === MC-GBO TABLE ===
\begin{table}[htbp]
    \centering
    \caption{Monte Carlo Gradient-Based Optimization (MC-GBO) hyperparameters.}
    \label{tab:params_mc_gbo}
    \begin{tabularx}{\textwidth}{l l X}
        \toprule
        \textbf{Parameter} & \textbf{Value} & \textbf{Description} \\
        \midrule
        $N_{MC}$ & $100$ & Total number of independent Monte Carlo random seeds evaluated. \\
        $N_{\text{parallel}}$ & $8$ & Number of optimizations executed concurrently. \\
        $T_{\max}$ & $180.0$\,s & Maximum allowable wall-clock computation time. \\
        \bottomrule
    \end{tabularx}
\end{table}

% === P-HO TABLE ===
\begin{longtable}{@{}p{0.24\textwidth}p{0.16\textwidth}p{0.54\textwidth}@{}}
    \caption{Probabilistic Homotopy Optimization (P-HO) hyperparameters.}
    \label{tab:params_pho}\\
        \toprule
        \textbf{Parameter} & \textbf{Value} & \textbf{Description} \\
        \midrule
\endfirsthead

        	oprule
        	extbf{Parameter} & \textbf{Value} & \textbf{Description} \\
        \midrule
\endhead

        \bottomrule
\endfoot

        $N_{\text{iters}}$ & $100$ & Maximum number of homotopy iterations. \\
        $T_{\max}$ & $180.0$\,s & Maximum allowable wall-clock computation time. \\
        $\rho_{\mathcal{A}}$ & $1.0$ & Solve-to-sample ratio threshold regulating phase transitions. \\
        $P_g$ & $0.3$ & Probability of explicitly targeting the goal parameter $\boldsymbol{\lambda} = \mathbf{1}_d$. \\
        $N_{\text{init}}$ & $20$ & Number of initial solutions generated at the relaxed state $\boldsymbol{\lambda} = \mathbf{0}_d$. \\
        $s$ & $0.9$ & Shrink factor for dynamic limits and exclusion zone constraints. \\
        Geodesic Seed & True & Initializes one baseline solve with the shortest-path geodesic slew. \\
        \bottomrule
\end{longtable}

% === RRT*-GBO TABLE ===
\begin{table}[htbp]
    \centering
    \caption{Algorithm-specific hyperparameters for Rapidly-exploring Random Tree Star with GBO (RRT*-GBO).}
    \label{tab:params_rrt_gbo}
    \begin{tabularx}{\textwidth}{l l X}
        \toprule
        \textbf{Parameter} & \textbf{Value} & \textbf{Description} \\
        \midrule
        $N_{\text{iters}}$ & $5000$ & Maximum number of tree expansion iterations. \\
        $T_{\max}$ & $180.0$\,s & Maximum allowable wall-clock computation time. \\
        $N_{\text{sol, max}}$ & $3$ & Maximum number of feasible solutions found before termination (only improved solutions increment). \\
        $r_{\text{prune}}$ & $0.1$\,rad & Radius threshold for pruning redundant nodes. \\
        $r_{\text{ball}}$ & $0.33$\,rad & Radius for neighborhood querying and rewiring consideration. \\
        $r_{\text{connect}}$ & $0.5\pi$\,rad & Goal connection threshold radius. \\
        $W_{\text{prune}}$ & $20$ & Pruning window size for tracking recent node addition attempts to dynamically adjust the pruning radius. \\
        $\varepsilon_{\text{cost}}$ & $10^{-4}$ & Minimum relative cost improvement tolerance for rewiring acceptance. \\
        Edge $n_{\text{via}}$ & $11$ & Number of via points used for discretizing local edge BVP solves. \\
        Edge $n_{\text{fcv}}$ & $3$ & Number of free interior control vertices per edge BVP solve. \\
        Truncation & True & Truncates out-of-bounds samples to the edge of $r_{\text{ball}}$ instead of rejecting. \\
        \bottomrule
    \end{tabularx}
\end{table}

% === IPSO TABLE ===
\begin{longtable}{@{}p{0.24\textwidth}p{0.16\textwidth}p{0.54\textwidth}@{}}
    \caption{Algorithm-specific hyperparameters for Inverse Dynamics Particle Swarm Optimization (IPSO).}
    \label{tab:params_ipso}\\
        \toprule
        \textbf{Parameter} & \textbf{Value} & \textbf{Description} \\
        \midrule
\endfirsthead

        \toprule
        \textbf{Parameter} & \textbf{Value} & \textbf{Description} \\
        \midrule
\endhead

        \bottomrule
\endfoot

        \multicolumn{3}{l}{\textit{Execution \& Swarm Dynamics}} \\
        $N_{\text{particles}}$ & $30$ & Number of particles in the swarm population. \\
        $N_{\text{runs}}$ & $20$ & Number of independent IPSO multistart runs. \\
        $N_{\text{parallel,IPSO}}$ & $4$ & Number of concurrent IPSO processes executed. \\
        $N_{\text{threads,IPSO}}$ & $2$ & threads allocated per IPSO process. \\
        $N_{\text{outer}}$ & $1000$ & Maximum number of outer loop iterations. \\
        $N_{\text{inner}}$ & $20$ & Number of internal optimization loops per outer cycle. \\
        $r_{\text{neighborhood}}$ & $3$ & Number of adjacent particles defining the local neighborhood. \\
        $v_{\text{fraction}}$ & $0.1$ & Search space fraction defining the maximum particle velocity clamping. \\
        Geodesic Seed & True & Initializes one particle with the shortest-path geodesic trajectory. \\
        \midrule
        \multicolumn{3}{l}{\textit{MRP Search-Space Limits}} \\
        $\theta_{\text{search}}$ & $5.0$\,rad & Fixed angular boundary applied component-wise to the MRP search space. \\
        $s_{\text{bound}}$ & $1.5$ & Margin multiplier applied to the geodesic magnitude (used if $\theta_{\text{search}}$ is null). \\
        \midrule
        \multicolumn{3}{l}{\textit{GBO Refinement}} \\
        GBO Refinement & True & Enables local gradient-based refinement on the final IPSO trajectories. \\
        Feasible Only & True & Restricts GBO refinement exclusively to strictly feasible IPSO outputs. \\
        $N_{\text{warmstarts}}$ & $20$ & Maximum number of top IPSO candidates passed to the GBO solver. \\
        $N_{\text{parallel,GBO}}$ & $8$ & Number of concurrent GBO refinement processes executed. \\
        \midrule
        \multicolumn{3}{l}{\textit{Weight Scheduling ($N_{\text{schedule}} = 300$ outer cycles)}} \\
        $r$ & $0.9$ & Static scalar multiplier for the combined inertia and cognitive velocity update. \\
        $w_0 \to w_f$ & $1.2 \to 0.6$ & Inertia weight linear decay schedule. \\
        $c_p$ & $1.5$ & Static cognitive weight for personal-best attraction. \\
        $c_{l0} \to c_{lf}$ & $2.0 \to 0.0$ & Local neighborhood weight linear decay schedule. \\
        $c_{g0} \to c_{gf}$ & $0.0 \to 2.0$ & Global swarm weight linear growth schedule. \\
        \midrule
        \multicolumn{3}{l}{\textit{Constraint Handling \& Convergence}} \\
        $\delta_0$ & $0.05$ & Initial additive keep-out tolerance for exclusion zones. \\
        $\gamma$ & $0.05$ & Adaptive decay factor for exclusion zone tolerance tightening. \\
        $N_{\delta}$ & $150$ & Cycle threshold triggering adaptive keep-out tolerance decay. \\
        $N_C$ & $50$ & Sliding history window size for evaluating cost convergence. \\
        $\varepsilon_{\text{tol}}$ & $10^{-6}$ & Relative cost improvement threshold required for convergence. \\
        \midrule
        \multicolumn{3}{l}{\textit{Penalties (Static Multipliers)}} \\
        $w_{\text{path}}, w_{\text{torque}}, w_{\omega}$ & $10.0$ & Respective penalty weights for EZ, torque, and velocity bounds. \\
        $w_{\text{violation}}$ & $10.0$ & Fixed penalty applied per active constraint violation count. \\
        \bottomrule
\end{longtable}

\section{Problem 2: 7-DOF Robotic Manipulator}
\label{sec:setup_robot}

\subsection{Physical Parameters}
\todoX{Include a table defining the \ac{MDH} parameters ($\alpha$, $a$, $\theta$, $d$) for all 7 links.}
\todoX{Include a table listing the mass, center of mass, and inertia tensor for each link.}
\todoX{List the specific joint position bounds $\mathbf{q}_{\max/\min}$, velocity bounds $\dot{\mathbf{q}}_{\max/\min}$, and torque bounds $\boldsymbol{\tau}_{\max/\min}$.}

A core component of the robotic manipulator model is its kinematic chain, which is defined using the Modified Denavit--Hartenberg (\ac{MDH}) parameters as described in Section~\ref{sec:manipulator-kinematics}. 
By chaining frame multiplication using the \ac{MDH} parameters and the fixed transformations connecting the inertial frame to the robot base and the last link to the end-effector, the homogeneous transformation from the inertial frame to the end-effector frame can be obtained.
Table~\ref{tab:MDH_parameters} lists the \ac{MDH} parameters for all links of the robotic manipulator, seven of which depend on the joint positions $q_1$ to $q_7$.

\begin{table}[htbp]
\centering
\caption{Modified DH parameters for the robotic manipulator kinematic chain.}
\begin{tabular}{c c c c c}
\hline
\label{tab:MDH_parameters}
$i$ &
$\alpha_{i-1}$ &
$a_{i-1}$ [mm] &
$\theta_{i-1}$ &
$d_{i-1}$ [mm] \\
\hline
1 & $0^\circ$    & 0 & $q_1$              & 310 \\
2 & $+90^\circ$  & 0 & $q_2$              & 0   \\
3 & $-90^\circ$  & 0 & $q_3 + 90^\circ$   & 400 \\
4 & $+90^\circ$  & 0 & $q_4$              & 0   \\
5 & $-90^\circ$  & 0 & $q_5 - 90^\circ$   & 390 \\
6 & $+90^\circ$  & 0 & $q_6 + 90^\circ$   & 0   \\
7 & $+90^\circ$  & 0 & $q_7 - 90^\circ$   & 0   \\
8 & $-90^\circ$ & 0 & $-180^\circ$       & 118 \\
\hline
\end{tabular}
\end{table}

The fixed transformations connecting the inertial frame to the robot base and frame 8 to the end-effector frame are
\begin{equation}
    \mathbf{T}_{i,0} = 
    \begin{bmatrix}
        -0.498791 & -0.866722 & 0 & 760.2 \\
        0.866722 & -0.498791 & 0 & -921.7 \\
        0 & 0 & 1 & 942.0 \\
        0 & 0 & 0 & 1
    \end{bmatrix},
    \label{eq:robot-base-transform}
\end{equation}
\begin{equation}
    \mathbf{T}_{\mathrm{8,ee}} = 
    \begin{bmatrix}
        -0.5411 & -0.2439 & 0.8048 & 35.0 \\
        -0.1669 & -0.9065 & -0.3875 & -30.0 \\
        0.8237 & -0.3440 & 0.4499 & 181.0 \\
        0 & 0 & 0 & 1
    \end{bmatrix},
    \label{eq:robot-tool-transform}
\end{equation}
with the translation components are given in millimeters.
The full conversion from the inertial frame to the end-effector frame is obtained by chaining the transformations as shown in \cref{eq:robot_forward_kinematics}.

The model dynamics were evaluated using a Maple-generated inverse dynamics model written in C. 
The model returns the mass matrix, centripetal-Coriolis vector, and gravity vector---used in  \cref{eq:robot_inverse_dynamics}---for the manipulator given the current joint positions and velocities.
The link physical parameters such as masses, centers of mass, and inertia tensors are stored within this generated dynamics model, rather than supplied explicitly at runtime.
The optimization uses this multi-rigid-body model directly and the link-specific physical parameters were not available in the project source.
As such, these parameters are not explicitly listed in this work.
However, the model does accept parameters for an added payload and joint motor inertias---which is used to calculate the cost function.
The payload parameters were $m_{load} = 2.725 \text{ kg}$ for the mass, ${}^{ee}\mathbf{p}_{load,com} = [0, 0, 0.1]^\top \text{m}$ for the center of mass position in the end-effector frame, and $\mathbf{J}_{load} = \mathbf{0}$ for the inertia matrix.
The joint motor inertias were specified as $\mathbf{J}_{\mathrm{m}} = \operatorname{diag}\left(4.121601, 4.121601, 1.384900, 1.384900, 0.541000, 0.600832, 0.600832\right)$ for joints 1 through 7.

Table~\cref{tab:robot_kindyn_limits} lists the joint position, velocity, and torque limits for the robotic manipulator on a per joint basis.

\begin{table}[htbp]
\centering
\caption{Robotic manipulator joint position, velocity, and torque limits.}
\begin{tabular}{c c c c c c c}
\hline
\label{tab:robot_kindyn_limits}
Joint &
$q_{\min}$ [rad] &
$q_{\max}$ [rad] &
$\dot{q}_{\min}$ [rad/s] &
$\dot{q}_{\max}$ [rad/s] &
$\tau_{\min}$ [N\,m] &
$\tau_{\max}$ [N\,m] \\
\hline
1 & -2.9671 & 2.9671 & -1.7453 & 1.7453 & -126 & 126 \\
2 & -2.0944 & 2.0944 & -1.7453 & 1.7453 & -76  & 76  \\
3 & -2.9671 & 2.9671 & -1.7453 & 1.7453 & -36  & 36  \\
4 & -2.0944 & 2.0944 & -1.7453 & 1.7453 & -36  & 36  \\
5 & -2.9671 & 2.9671 & -2.6180 & 2.6180 & -11  & 11  \\
6 & -0.7854 & 1.3963 & -1.7453 & 1.7453 & -16  & 16  \\
7 & -0.7854 & 2.3562 & -1.7453 & 1.7453 & -16  & 16  \\
\hline
\end{tabular}
\end{table}

\subsection{Target (Ball) Parameters}
\todoX{Specify the ball's initial position ${}^i\mathbf{p}_{b,0}$, initial velocity ${}^i\mathbf{v}_{b,0}$, and gravity vector ${}^i\mathbf{g}$.}

\subsection{B-Spline Discretization}
\todoX{Specify B-spline degree $p=3$ and control vertex count $n_{ctrl}$.}
\todoX{State the catch time horizon bounds $t_f \in [0.06, 1.8]$ s.}
\todoX{State the oversampling ratio ($n_{via} \approx 2 n_{ctrl}$).}

\subsection{Test Scenarios}
\todoX{Define the initial robot configuration $\mathbf{q}_0$.}
\todoX{Detail any variations in ball throwing arcs (e.g., to force different kinematic homotopy classes).}

\subsection{Algorithm-Specific Hyperparameters}
\todoX{RRT*-GBO: Joint-space distance metrics for $r_{ball,joint}$ and task-space metrics for $r_{connect}$.}
\todoX{P-HO: Initial $\pm 2\pi$ joint position relaxations, 10x dynamic limit relaxations, and initial guess generation strategies.}